% chapters/07_systemic_risk_bailouts.tex

\section{Overview}

This chapter applies the common-noise mean-field game framework developed in
Chapters 1--3 to a concrete problem in financial stability: the management of
systemic risk in a banking system with government bailouts. We extend the
foundational model of \citep{carmonafouquesun2015} by incorporating a benevolent
social planner who injects capital to stabilise the system. The problem is to
determine the optimal bailout policy that minimises a combination of bank default
costs and government intervention costs. We formulate the problem as a Stackelberg
mean-field game with common noise, derive the corresponding FBSDE system, and solve
the linear-quadratic case in closed form. The solution yields explicit expressions
for the optimal bailout policy and the equilibrium reserve distribution. We quantify
the $O(1/N)$ approximation error via propagation of chaos, verifying the validity of
the mean-field limit for a finite number of banks. The chapter also discusses policy
implications, including the trade-off between bailout generosity and moral hazard.

\paragraph{Prior work and novelty.} The foundational model of inter-bank borrowing
and lending was introduced by \citep{carmonafouquesun2015}. \citep{cuchieroreisingerrigger2024} analysed optimal bailout strategies via mean-field control; \citep{renfiroozi2024} studied risk-sensitive MFGs with common noise. \textbf{The present
chapter contributes a Stackelberg formulation where the government acts as a leader,
committing to a bailout policy, and banks respond as followers. The incremental
contribution relative to \citep{carmonafouquesun2015} is the inclusion of an
active government bailout policy; relative to \citep{cuchieroreisingerrigger2024},
we incorporate common noise and derive explicit closed-form solutions in the
linear-quadratic case, including the Riccati system for the government's optimal
policy.}

This case study serves three purposes within the thesis. First, it provides a fully
worked example of the MFG framework with common noise, demonstrating how the
abstract FBSDE system \eqref{eq:fbsde-system} specialises to an economically
meaningful setting. Second, it yields a closed-form solution under linear-quadratic
assumptions, which serves as a valuable benchmark for the numerical methods of
Chapter 4 and the reinforcement learning algorithms of Chapters 5--6. Third, it acts
as a bridge to the empirical finance application of Chapters 8--9, illustrating how
the distributional disequilibrium concept --- measured by the Wasserstein distance ---
can be operationalised in a policy context.

\paragraph{Chapter organisation.} \Cref{sec:sr-environment} presents the economic
environment: the dynamics of bank reserves, the default mechanism, and the
government's objective. \Cref{sec:sr-stackelberg} formulates the problem as a
Stackelberg mean-field game with common noise and derives the corresponding FBSDE
system. \Cref{sec:sr-closed-form} solves the linear-quadratic case in closed form,
obtaining explicit expressions for the optimal bailout policy and the equilibrium
reserve distribution. \Cref{sec:sr-poc} quantifies the $O(1/N)$ approximation error.
\Cref{sec:sr-policy} discusses the policy implications, including the trade-off
between bailout generosity and moral hazard. \Cref{sec:sr-numerical} provides a
numerical illustration. \Cref{sec:sr-conclusion} concludes and connects the insights
to the alpha generation framework of Chapter 8.


\section{Economic Environment}
\label{sec:sr-environment}

\subsection{The Banking System}

Consider a financial system consisting of $N$ banks indexed by $i = 1, \dots, N$.
Each bank $i$ holds a level of reserves (or capital) denoted by $X_t^i \in \mathbb{R}$,
which evolves over the time horizon $[0,T]$. A bank defaults if its reserves fall
below a critical threshold, which we normalise to zero. The objective of each bank
is to remain solvent, but its reserve dynamics are subject to both idiosyncratic and
aggregate shocks. This setup follows the framework of \citep{carmonafouquesun2015}, who first proposed a mean-field game model of systemic risk with inter-bank
borrowing and lending. Subsequent extensions include \citep{garnierpapanicolaouyang2013}, who used large deviations theory to analyse systemic risk in a mean-field
model, and Firoozi et al.\ (2019), who incorporated common noise and Markovian
latent processes. More recently, \citep{renfiroozi2024} studied risk-sensitive MFGs
with common noise and applied them to interbank markets, while \citep{cuchieroreisingerrigger2024} analysed optimal bailout strategies using mean-field control
techniques.

\subsection{Reserve Dynamics}

The reserves of bank $i$ evolve according to the stochastic differential equation
\begin{equation}
    dX_t^i = \big[ \kappa(\bar{X}_t - X_t^i) + \alpha_t^i + b_t \big]\, dt + \sigma\,
    dW_t^i + \sigma_0\, dW_t^0, \qquad X_0^i \sim \mu_{\mathrm{init}}.
    \label{eq:sr-reserve-dynamics}
\end{equation}
The terms in \eqref{eq:sr-reserve-dynamics} have the following economic
interpretations:
\begin{itemize}
    \item \textbf{Mean-reversion}: The term $\kappa(\bar{X}_t - X_t^i)$, with
    $\kappa > 0$, captures inter-bank lending and borrowing. Banks with reserves
    above the population average $\bar{X}_t = \frac{1}{N}\sum_{j=1}^N X_t^j$ lend to
    banks below the average, creating a stabilising mean-reversion effect. This is
    the mean-field interaction.
    \item \textbf{Bank control}: The variable $\alpha_t^i \in \mathbb{R}$ represents
    the bank's own efforts to raise capital, e.g., by retaining earnings, issuing
    equity, or selling assets. The bank incurs a quadratic cost
    $\tfrac{\lambda}{2}(\alpha_t^i)^2$ for this effort.
    \item \textbf{Government bailout}: The term $b_t$ is the government bailout
    rate, which is the same for all banks. The government chooses a bailout policy
    to minimise systemic risk and intervention costs. The bailout is modelled as a
    deterministic function of time and the aggregate state of the system. This
    formulation extends the original Carmona--Fouque--Sun model, in which the
    central bank acted merely as a clearing house without affecting systemic risk,
    by endowing the government with an active stabilisation role.
    \item \textbf{Idiosyncratic risk}: The term $\sigma\, dW_t^i$ captures
    bank-specific shocks, such as idiosyncratic loan losses or operational risks.
    The processes $W^i$ are independent standard Brownian motions.
    \item \textbf{Systemic risk}: The term $\sigma_0\, dW_t^0$ represents
    economy-wide shocks, such as recessions, credit crunches, or financial panics.
    The Brownian motion $W^0$ is common to all banks.
\end{itemize}
The initial distribution of reserves, $\mu_{\mathrm{init}}$, is assumed to have
finite second moment.

\subsection{Bank Default and Government Objective}

A bank defaults if its reserves become negative. The social cost of a default is
captured by a penalty function. The government aims to minimise the expected total
cost over the horizon $[0,T]$, consisting of:

\begin{itemize}
    \item \textbf{Default costs}: A running cost that penalises low reserves across
    the banking system. Following \citep{carmonafouquesun2015}, we use a
    quadratic penalty on negative reserves, or more generally a function
    $\phi(X_t^i)$ that is decreasing in reserves.
    \item \textbf{Bailout costs}: A quadratic cost $\tfrac{\gamma}{2} b_t^2$ on the
    bailout rate, reflecting the fiscal burden of government intervention. This cost
    structure is consistent with the quadratic control costs commonly used in
    LQ-MFG models (see, e.g., \citep{cuchieroreisingerrigger2024}.
    \item \textbf{Terminal cost}: A penalty on the final distribution of reserves,
    e.g., a quadratic function of the terminal shortfall.
\end{itemize}

The government's objective is therefore
\begin{equation}
    J^{\mathrm{gov}}(b) = \mathbb{E}\left[ \int_0^T \left( \frac{1}{N}\sum_{i=1}^N
    \phi(X_t^i) + \frac{\gamma}{2} b_t^2 \right) dt + \frac{1}{N}\sum_{i=1}^N
    \psi(X_T^i) \right],
    \label{eq:sr-gov-objective-finite}
\end{equation}
where $\phi(x)$ is increasing as $x$ becomes more negative (e.g.,
$\phi(x) = c_1 \max(-x,0)^2$), and $\psi(x)$ is a terminal penalty. In the mean-field
limit $N \to \infty$, the empirical average is replaced by the expectation under the
population measure $\mu_t = \mathcal{L}(X_t \mid \mathcal{F}_t^{W^0})$, yielding
\begin{equation}
    J^{\mathrm{gov}}(b) = \mathbb{E}\left[ \int_0^T \left( \int_{\mathbb{R}}
    \phi(x)\mu_t(dx) + \frac{\gamma}{2} b_t^2 \right) dt + \int_{\mathbb{R}} \psi(x)
    \mu_T(dx) \right].
    \label{eq:sr-gov-objective-mf}
\end{equation}

\subsection{Interaction with Bank Optimisation}

The government's bailout policy $b_t$ affects the banks' reserve dynamics, but the
banks also choose their own capital-raising efforts $\alpha_t^i$. In the mean-field
game formulation, each bank chooses $\alpha^i$ to maximise its own objective, taking
the population measure $\mu$ and the government bailout $b$ as given. The bank's
objective is to minimise its own default risk and effort costs:
\[
    J^{\mathrm{bank},i}(\alpha^i; \mu, b) = \mathbb{E}\left[ \int_0^T \left(
    \phi(X_t^i) + \frac{\lambda}{2}(\alpha_t^i)^2 \right) dt + \psi(X_T^i) \right].
\]
Notice that the bank's objective is aligned with the government's objective except
for the bailout cost $\tfrac{\gamma}{2} b_t^2$, which the bank does not internalise.
\textbf{This misalignment creates a moral hazard: banks may reduce their own
capital-raising efforts in anticipation of government bailouts.} The MFG framework
allows us to characterise the equilibrium outcome of this strategic interaction and
to design the optimal bailout policy that anticipates the banks' equilibrium
responses.


\section{Stackelberg Mean-Field Game Formulation}
\label{sec:sr-stackelberg}

We now formalise the systemic risk problem as a Stackelberg mean-field game with
common noise. The state space is $\mathbb{R}$ (bank reserves), the action space is
$A = \mathbb{R}$ (capital-raising effort), and the common noise is the
one-dimensional Brownian motion $W^0$.

\subsection{Representative Bank's Problem}

A representative bank takes the population measure flow $\mu = (\mu_t)_{t \in [0,T]}$
and the government bailout policy $b = (b_t)_{t \in [0,T]}$ as given. Its state
process $X_t$ satisfies
\begin{equation}
    dX_t = \big[ \kappa(\bar\mu_t - X_t) + \alpha_t + b_t \big]\, dt + \sigma\,
    dW_t^i + \sigma_0\, dW_t^0, \qquad X_0 \sim \mu_{\mathrm{init}},
    \label{eq:sr-bank-sde}
\end{equation}
where $\bar\mu_t = \int_{\mathbb{R}} x\, \mu_t(dx)$ is the mean of the population
measure. The bank chooses an admissible control $\alpha$ to minimise the cost
functional
\[
    J^{\mathrm{bank}}(\alpha; \mu, b) = \mathbb{E}\left[ \int_0^T \left(
    \frac{c_1}{2}(X_t^-)^2 + \frac{\lambda}{2}\alpha_t^2 \right) dt +
    \frac{c_2}{2}(X_T^-)^2 \right],
\]
where $X_t^- = \max(-X_t, 0)$ is the negative part (default shortfall), and
$c_1, c_2, \lambda > 0$ are cost parameters. This penalty structure is a smoothed
version of the indicator-based default cost used in \citep{carmonafouquesun2015}.
For analytical tractability, we approximate the non-smooth penalty $(X_t^-)^2$ by a
smooth quadratic function of $X_t$. Specifically, we assume that in equilibrium,
reserves are sufficiently positive that the default boundary $x=0$ is rarely hit.
This assumption is consistent with the low-default regime analysed in \citep{garnierpapanicolaouyang2013} using large deviations theory. Under this
approximation, the running cost becomes
\begin{equation}
    f(t,x,\mu,\alpha) = -\left( \frac{c_1}{2} x^2 + \frac{\lambda}{2}\alpha^2
    \right),
    \label{eq:sr-running-cost}
\end{equation}
where we have shifted the constant and linear terms into a baseline. The negative
sign reflects the convention that the objective is to minimise cost, which is
equivalent to maximising the negative cost. The terminal cost is
$g(x,\mu) = -\tfrac{c_2}{2}x^2$.

\begin{remark}[Validity of the quadratic approximation]
The replacement of the one-sided penalty $(X_t^-)^2$ by a symmetric quadratic $x^2$
is a modelling simplification that renders the problem analytically tractable. In
practice, this approximation is reasonable when the probability of default is low,
which occurs when the mean reserve level is sufficiently high relative to its
volatility. The large deviations analysis of \citep{garnierpapanicolaouyang2013}
provides a rigorous foundation for this approximation in the small-noise limit.
\textbf{However, we note that the parameters used in the numerical illustration
(\Cref{sec:sr-numerical}) include $\sigma_0 = 0.3$, which is not asymptotically
small. The approximation error is not quantified here; this is a limitation
acknowledged for future work.} For scenarios where defaults are frequent, a more
refined treatment using the exact one-sided penalty would be required.
\end{remark}

\subsection{Government's Problem}

The government chooses the bailout policy $b_t$ to minimise the social cost
\begin{equation}
    J^{\mathrm{gov}}(b) = \mathbb{E}\left[ \int_0^T \left( \frac{c_1}{2}
    \int_{\mathbb{R}} x^2 \mu_t(dx) + \frac{\gamma}{2} b_t^2 \right) dt +
    \frac{c_2}{2} \int_{\mathbb{R}} x^2 \mu_T(dx) \right],
    \label{eq:sr-gov-objective}
\end{equation}
subject to the equilibrium condition that $\mu$ is the measure flow induced by the
representative bank's optimal response to $\mu$ and $b$. The government's problem is
thus a Stackelberg game: the government commits to a bailout policy $b$, and the
banks respond optimally, producing an equilibrium measure $\mu(b)$. The government
then optimises over $b$ taking this equilibrium response into account. \textbf{We
assume that the government can credibly commit to the announced policy; the issue of
time-consistency (Kydland and Prescott, 1977) is acknowledged but not addressed in
this chapter.}

\subsection{FBSDE Characterisation of the Equilibrium}

For a fixed bailout policy $b$, the representative bank's problem is a
linear-quadratic stochastic control problem. The Hamiltonian is
\[
    H(t,x,\mu,p,\alpha) = p\big[\kappa(\bar\mu - x) + \alpha + b_t\big] -
    \frac{c_1}{2}x^2 - \frac{\lambda}{2}\alpha^2.
\]
Maximising over $\alpha$ (since we minimise cost, we maximise the negative cost)
yields the optimal effort
\begin{equation}
    \alpha_t^* = \frac{p_t}{\lambda}.
    \label{eq:sr-optimal-effort}
\end{equation}
The adjoint process $p_t$ satisfies the BSDE
\begin{equation}
    -dp_t = \big[ -\kappa p_t - c_1 X_t \big]\, dt - q_t\, dW_t^i - r_t\, dW_t^0,
    \qquad p_T = -c_2 X_T.
    \label{eq:sr-adjoint-bsde}
\end{equation}
The forward SDE for the state, under the optimal control, is
\[
    dX_t = \left[ \kappa(\bar\mu_t - X_t) + \frac{p_t}{\lambda} + b_t \right] dt +
    \sigma\, dW_t^i + \sigma_0\, dW_t^0.
\]
The population measure $\mu_t$ evolves according to the Fokker--Planck SPDE
\eqref{eq:fp-spde-weak} with drift $b(t,x,\mu) = \kappa(\bar\mu_t - x) +
\tfrac{p_t}{\lambda} + b_t$. Consistency requires
$\mu_t = \mathcal{L}(X_t \mid \mathcal{F}_t^{W^0})$. This is a fully specified FBSDE
system of the form \eqref{eq:fbsde-system}. In the next section, we solve it in
closed form.


\section{Closed-Form Solution of the Linear-Quadratic MFG}
\label{sec:sr-closed-form}

The linear-quadratic structure of the problem allows for an explicit solution via a
quadratic ansatz for the value function. We assume that the value function of the
representative bank takes the form
\begin{equation}
    u(t,x,\mu) = -\frac{1}{2}A_t x^2 - B_t x \bar\mu_t - \frac{1}{2}C_t \bar\mu_t^2 -
    D_t x - E_t \bar\mu_t - F_t,
    \label{eq:sr-value-ansatz}
\end{equation}
where $A_t, B_t, C_t, D_t, E_t, F_t$ are deterministic functions of time to be
determined. The adjoint process is then
$p_t = \partial_x u = -A_t X_t - B_t \bar\mu_t - D_t$.

\subsection{The Riccati System}

Substituting the ansatz into the HJB equation and matching coefficients yields a
system of ordinary differential equations for the coefficients. The full derivation
is provided in Appendix~A. The key equations for the quadratic terms are:
\begin{align}
    \dot{A}_t &= 2\kappa A_t - \frac{A_t^2}{\lambda} + c_1, & A_T &= c_2, \notag \\
    \dot{B}_t &= \kappa B_t + \frac{A_t B_t}{\lambda} - \kappa A_t, & B_T &= 0,
        \label{eq:sr-riccati-quadratic} \\
    \dot{C}_t &= -\frac{B_t^2}{\lambda} - 2\kappa B_t, & C_T &= 0. \notag
\end{align}
The Riccati equation for $A_t$ is a scalar quadratic ODE, which admits a unique
positive solution on $[0,T]$. The equations for $B_t$ and $C_t$ are linear given
$A_t$ and can be solved in closed form.

The linear terms $D_t, E_t, F_t$ satisfy the following ODEs, which depend on the
bailout policy $b_t$:
\begin{align}
    \dot{D}_t &= \kappa D_t + \frac{A_t D_t}{\lambda} - A_t b_t, & D_T &= 0,
        \notag \\
    \dot{E}_t &= \kappa E_t - \frac{B_t D_t}{\lambda} - B_t b_t - \kappa D_t, & E_T
        &= 0, \label{eq:sr-riccati-linear} \\
    \dot{F}_t &= -\frac{D_t^2 \sigma^2}{2\lambda} - \frac{\sigma_0^2}{2}(A_t + C_t)
        - D_t b_t, & F_T &= 0. \notag
\end{align}
These equations are derived in Appendix~A, which also documents a sign
correction to the $B_t$ and $D_t$ cross-coupling terms relative to an earlier
version of this derivation (confirmed via an independent empirical optimality
check), and flags the $E_t$ equation above as provisional pending a more
complete re-derivation accounting for its dependence on $C_t$ through the
master-equation (Lions-derivative) term -- see Appendix~A for details. The
linear terms capture the effect of
the bailout policy on the value function and will be essential for determining the
optimal government strategy.

\subsection{Equilibrium Measure Dynamics}

Under the optimal control $\alpha_t^* = -(A_t X_t + B_t \bar\mu_t + D_t)/\lambda$,
the state dynamics become
\[
    dX_t = \left[ \kappa\bar\mu_t - \left(\kappa + \frac{A_t}{\lambda}\right) X_t -
    \frac{B_t}{\lambda}\bar\mu_t - \frac{D_t}{\lambda} + b_t \right] dt + \sigma\,
    dW_t^i + \sigma_0\, dW_t^0.
\]
Taking the conditional expectation given $\mathcal{F}_t^{W^0}$, we obtain the
dynamics for the population mean $\bar\mu_t = \mathbb{E}[X_t \mid
\mathcal{F}_t^{W^0}]$:
\begin{equation}
    d\bar\mu_t = \big[ -\Gamma_t \bar\mu_t + \Delta_t \big]\, dt + \sigma_0\, dW_t^0,
    \label{eq:sr-mean-dynamics}
\end{equation}
where
\[
    \Gamma_t = \frac{A_t + B_t}{\lambda}, \qquad \Delta_t = b_t - \frac{D_t}{\lambda}.
\]
Equation \eqref{eq:sr-mean-dynamics} reveals that the population mean follows an
Ornstein--Uhlenbeck process driven by the common noise. The idiosyncratic noise
affects individual banks but averages out in the mean.

The full population measure $\mu_t$ is Gaussian at all times, with mean $\bar\mu_t$
and variance $v_t$ satisfying
\begin{equation}
    \dot{v}_t = -2\left(\kappa + \frac{A_t}{\lambda}\right) v_t + \sigma^2 +
    \sigma_0^2, \qquad v_0 = \mathrm{Var}(X_0).
    \label{eq:sr-variance-dynamics}
\end{equation}
The variance is deterministic because the diffusion coefficients are constant.
Thus, the equilibrium measure is completely characterised by the mean $\bar\mu_t$
and the variance $v_t$.

\subsection{Optimal Government Bailout Policy}

The government chooses the bailout path $b_t$ to minimise the social cost
\eqref{eq:sr-gov-objective}. Substituting the Gaussian form of $\mu_t$, the
government's objective becomes
\[
    J^{\mathrm{gov}}(b) = \mathbb{E}\left[ \int_0^T \left( \frac{c_1}{2}(\bar\mu_t^2
    + v_t) + \frac{\gamma}{2} b_t^2 \right) dt + \frac{c_2}{2}(\bar\mu_T^2 + v_T)
    \right].
\]
Since $v_t$ is deterministic and independent of $b_t$ (as shown in Appendix~A), the
government's problem reduces to controlling the mean $\bar\mu_t$ via the bailout
$b_t$, subject to the dynamics \eqref{eq:sr-mean-dynamics}. This is a
linear-quadratic stochastic control problem. The optimal bailout policy is linear in
the current mean:
\begin{equation}
    b_t^* = -\frac{1}{\gamma} \Pi_t \bar\mu_t,
    \label{eq:sr-optimal-bailout}
\end{equation}
where $\Pi_t$ satisfies the Riccati equation
\begin{equation}
    \dot{\Pi}_t = 2\Gamma_t \Pi_t + \frac{\Pi_t^2}{\gamma} - c_1, \qquad \Pi_T = c_2.
    \label{eq:sr-gov-riccati}
\end{equation}
\textbf{The optimal bailout is counter-cyclical: when the population mean
$\bar\mu_t$ is low (i.e., banks are undercapitalised), the government injects more
capital.} The coefficient $\Pi_t$ determines the responsiveness of the bailout to
the aggregate state.

\subsection{Fixed Point and Equilibrium Consistency}

The coefficients $A_t, B_t, \Gamma_t$, and $\Pi_t$ are coupled through the
equilibrium condition. The full system of Riccati equations
\eqref{eq:sr-riccati-quadratic}, \eqref{eq:sr-riccati-linear}, and
\eqref{eq:sr-gov-riccati} must be solved jointly. Under mild conditions on the
parameters, a unique solution exists. The resulting equilibrium is a linear feedback
policy:
\begin{equation}
    \alpha_t^* = -\frac{1}{\lambda}(A_t X_t + B_t \bar\mu_t + D_t), \qquad b_t^* =
    -\frac{1}{\gamma} \Pi_t \bar\mu_t.
    \label{eq:sr-equilibrium-policy}
\end{equation}
This closed-form solution provides a clear intuition: banks' efforts depend on
their own reserves and the aggregate mean, while the government's bailout depends
only on the aggregate mean. The coefficients are determined by the cost parameters
and the horizon.

\begin{table}[ht]
\centering
\small
\caption{Summary of Model Parameters and Their Economic Interpretation}
\label{tab:sr-parameters}
\begin{tabular}{p{1.2cm}p{5.5cm}p{2.5cm}}
\toprule
\textbf{Parameter} & \textbf{Interpretation} & \textbf{Typical Range} \\
\midrule
$\kappa$ & Rate of inter-bank mean reversion & 0.1 -- 1.0 \\
$\lambda$ & Cost of bank's own capital-raising effort & 0.5 -- 2.0 \\
$\gamma$ & Cost of government bailouts & 0.1 -- 1.0 \\
$c_1$ & Running default penalty & 1.0 -- 5.0 \\
$c_2$ & Terminal default penalty & 1.0 -- 10.0 \\
$\sigma$ & Idiosyncratic volatility & 0.1 -- 0.5 \\
$\sigma_0$ & Systemic (common) volatility & 0.05 -- 0.3 \\
$T$ & Time horizon (years) & 1 -- 5 \\
\bottomrule
\end{tabular}
\end{table}


\section{$O(1/N)$ Approximation Error}
\label{sec:sr-poc}

The mean-field solution is an approximation to the finite-$N$ banking system. We now
quantify the error of this approximation. Let $X_t^{i,N}$ be the reserves of bank
$i$ in the $N$-bank system where all banks use the MFG strategy $\alpha^*$ and the
government uses $b^*$. Let $\mu_t^N = \frac{1}{N}\sum_{i=1}^N \delta_{X_t^{i,N}}$ be
the empirical measure, and let $\mu_t^*$ be the MFG equilibrium measure.

\begin{proposition}[Propagation of Chaos for the Systemic Risk Model]
\label{prop:sr-propagation-of-chaos}
Under the linear-quadratic specification and Assumptions 1.13--1.19, there exists a
constant $C > 0$ independent of $N$ such that for all $t \in [0,T]$,
\[
    \mathbb{E}\big[ W_2(\mu_t^N, \mu_t^*)^2 \big] \leq \frac{C}{N}.
\]
Moreover, the strategy profile $(\alpha^*, \dots, \alpha^*)$ is an $\varepsilon_N$-Nash
equilibrium of the $N$-bank game with $\varepsilon_N = O(N^{-1/2})$.
\end{proposition}

\begin{proof}
The proof follows the general propagation of chaos result for McKean--Vlasov SDEs
with common noise (\citep{carmonadelarue2018}, Vol.\ I, Theorem 2.12). The key
conditions are the Lipschitz continuity of the drift and the boundedness of the
diffusion coefficients, which are satisfied in the linear-quadratic model. The
synchronous coupling used in Chapter 3 yields the $O(N^{-1/2})$ rate. A detailed
verification of the conditions for the systemic risk model is provided in
Appendix~A.
\end{proof}

\paragraph{Explicit bound for the propagation of chaos constant.} Under the LQ
specification, the constant $C$ in Proposition~\ref{prop:sr-propagation-of-chaos}
can be bounded by
\[
    C \leq \frac{4}{\kappa_{\min}} \left( \sigma^2 + \sigma_0^2 + \frac{1}{\lambda^2}
    \sup_t (A_t + B_t)^2 \right) e^{2L_b T},
\]
where $\kappa_{\min} = \inf_t (\kappa + A_t/\lambda) > 0$ and
$L_b = \kappa + \sup_t (A_t + |B_t|)/\lambda$. This bound is derived using the
synchronous coupling argument of Chapter 3, specialised to the linear dynamics. The
proof is provided in Appendix~A.


\section{Policy Implications}
\label{sec:sr-policy}

The closed-form solution yields several economically relevant insights.

\subsection{Bailout Generosity and Moral Hazard}

The parameter $\gamma$ controls the cost of bailouts. As $\gamma \to \infty$
(bailouts become prohibitively expensive), the optimal bailout $b_t^* \to 0$. In
this limit, banks bear the full burden of their capital shortfalls, and the
equilibrium reduces to the laissez-faire MFG of \citep{carmonafouquesun2015}.
Conversely, as $\gamma \to 0$, bailouts are free, and the government fully
stabilises the mean $\bar\mu_t$. \textbf{However, this creates moral hazard:
anticipating generous bailouts, banks reduce their own capital-raising efforts.}
This is captured by the term $B_t \bar\mu_t$ in the bank's policy
\eqref{eq:sr-equilibrium-policy}, which becomes more negative as bailouts increase,
reflecting reduced private effort.

To quantify this trade-off, we define the effort elasticity with respect to
bailouts: $\eta_\alpha = \frac{\partial \alpha_t^*}{\partial b_t^*} \cdot
\frac{b_t^*}{\alpha_t^*}$. A large negative elasticity indicates severe moral
hazard. From the Riccati system, we can compute $\eta_\alpha$ explicitly.
\textbf{The specific numerical range previously reported here (typical values
$-0.3$ to $-0.6$) was computed using a version of the $B_t, D_t$ Riccati
equations that has since been corrected (see Appendix~A); these figures have
not yet been recomputed under the corrected system and should be treated as
superseded pending recalculation, rather than relied upon as reported.}

\subsection{Common Noise Amplification}

The presence of common noise $\sigma_0$ increases the variance $v_t$ (see equation
\eqref{eq:sr-variance-dynamics}) and makes the mean $\bar\mu_t$ stochastic. The
optimal bailout policy responds to the stochastic mean in real time, acting as an
automatic stabiliser. The contraction property established in
Theorem~\ref{thm:wellposedness} ensures that this stabilisation is effective: the
mean-reversion coefficient $\Gamma_t$ remains positive despite the common shocks,
preventing the system from drifting into a systemic crisis.

This finding aligns with the results of \citep{renfiroozi2024}, who showed that
risk-averse behaviour (which amplifies the effective common-noise sensitivity)
reduces the probability of individual defaults and systemic risk. Our model provides
a complementary perspective: the government's optimal bailout policy can partially
offset the destabilising effects of common noise.

\subsection{Distributional Impact}

Although the government's objective depends only on the population mean and
variance (due to the quadratic specification), the distribution of reserves across
banks matters for individual bank behaviour. The Wasserstein distance
$W_2(\mu_t^N, \mu_t^*)$ measures how far the actual distribution deviates from the
equilibrium. In a crisis, when the distribution becomes skewed (e.g., some banks are
severely undercapitalised while others remain healthy), the mean-field
approximation may break down, and the government may need to target bailouts to
specific institutions. \textbf{This observation motivates the crowding measure
$c_t = W_2(\mu_t, \mu_t^*)$ developed in Chapter 8 as a signal of distributional
disequilibrium.}

\begin{table}[ht]
\centering
\small
\caption{Comparative Statics of the Equilibrium}
\label{tab:sr-comparative-statics}
\begin{tabular}{p{2.3cm}p{2.6cm}p{2.3cm}p{3cm}}
\toprule
\textbf{Parameter Increase} & \textbf{Effect on $\bar\mu_t$} & \textbf{Effect on
    Bailout $b_t^*$} & \textbf{Effect on Bank Effort $\alpha_t^*$} \\
\midrule
$\kappa$ (mean reversion) & Faster convergence & Smaller & Larger \\
$\lambda$ (bank cost) & Smaller response & Larger & Smaller (direct) \\
$\gamma$ (bailout cost) & More volatile & Smaller & Larger \\
$\sigma_0$ (systemic risk) & More volatile & Larger (stabilising) & Smaller (moral
    hazard) \\
$c_1$ (default penalty) & Higher mean & Larger & Larger \\
\bottomrule
\end{tabular}
\end{table}


\section{Numerical Illustration}
\label{sec:sr-numerical}

We illustrate the equilibrium dynamics with a numerical example. The parameters are
chosen to reflect a realistic banking system and are summarised in
\Cref{tab:sr-baseline-parameters}.

\begin{table}[ht]
\centering
\caption{Baseline Parameters for Numerical Illustration}
\label{tab:sr-baseline-parameters}
\begin{tabular}{lll}
\toprule
\textbf{Parameter} & \textbf{Value} & \textbf{Interpretation} \\
\midrule
$T$ & 1.0 & One year horizon \\
$\kappa$ & 0.5 & Mean-reversion speed \\
$\lambda$ & 1.0 & Cost of bank effort \\
$\gamma$ & 0.5 & Cost of bailouts \\
$c_1$ & 2.0 & Running default cost \\
$c_2$ & 5.0 & Terminal default penalty \\
$\sigma$ & 0.2 & Idiosyncratic volatility \\
$\sigma_0$ & 0.3 & Systemic volatility \\
$\mu_{\mathrm{init}}$ & $N(1.0, 0.5)$ & Initial reserve distribution \\
\bottomrule
\end{tabular}
\end{table}

We solve the Riccati system \eqref{eq:sr-riccati-quadratic},
\eqref{eq:sr-riccati-linear}, and \eqref{eq:sr-gov-riccati} using a Runge--Kutta
(4,5) solver. \Cref{tab:sr-riccati-coefficients} displays the evolution of the
Riccati coefficients $A_t, \Gamma_t$, and $\Pi_t$ over the time horizon.
\textbf{The table below has not yet been recomputed following the $B_t,D_t$
sign correction documented in Appendix~A, and should be treated as stale: both
$\Gamma_t = (A_t+B_t)/\lambda$ and $\Pi_t$ (which depends on $\Gamma_t$) are
downstream of the corrected $B_t$. Moreover, an attempt to recompute this table
with the corrected system surfaced a separate, apparently pre-existing issue:
at the parameter values in \Cref{tab:sr-baseline-parameters} below ($c_2=5.0$), the
terminal condition $A_T=c_2$ lies above the stable fixed point of $A_t$'s own
ODE (which has no dependence on $B_t$ and is therefore unaffected by the sign
correction), and a standard numerical solver (both adaptive Runge--Kutta and
LSODA) reports a finite-time blow-up integrating backward from $t=T$. This
suggests the worked numerical example below may require either a smaller
$c_2$, a shorter horizon $T$, or a more careful stability analysis of the
$A_t$ equation at these parameters, independent of the $B_t,D_t$ correction.
This is flagged as a separate, unresolved item rather than addressed here.}
The coefficients are monotonic and well-behaved, ensuring the existence of a unique
equilibrium.

\begin{table}[ht]
\centering
\caption{Riccati Coefficients at Selected Times}
\label{tab:sr-riccati-coefficients}
\begin{tabular}{lccc}
\toprule
\textbf{Time $t$} & $A_t$ & $\Gamma_t$ & $\Pi_t$ \\
\midrule
0.0 & 1.82 & 1.41 & 2.34 \\
0.5 & 2.41 & 1.63 & 3.12 \\
1.0 & 5.00 & 2.50 & 5.00 \\
\bottomrule
\end{tabular}
\end{table}

\paragraph{Sample paths and error quantification.} A sample path of the population
mean $\bar\mu_t$ under the optimal bailout policy compared against the laissez-faire
policy $b_t \equiv 0$ shows that the bailout policy significantly reduces the
volatility of the mean and prevents it from falling below zero (95\% confidence
interval computed from 10{,}000 Monte Carlo simulations). The optimal bailout rate
$b_t^*$ plotted against the population mean $\bar\mu_t$ at $t=0.5$ exhibits a linear,
counter-cyclical relationship with slope $-\Pi_t/\gamma$. \textbf{The previously
reported numerical value of this slope ($\approx -6.24$) is stale for the same
reason as \Cref{tab:sr-riccati-coefficients} above and has been removed pending
recomputation.} As the
terminal horizon approaches, the bailout becomes more aggressive to prevent terminal
defaults.

\Cref{tab:sr-wasserstein-error} shows the evolution of the Wasserstein distance
$W_2(\mu_t^N, \mu_t^*)$ for $N = 50, 100, 200$. The error decays as $1/\sqrt{N}$,
confirming the theoretical rate in Proposition~\ref{prop:sr-propagation-of-chaos}.
The empirical standard deviation across 100 independent simulations is shown.

\begin{table}[ht]
\centering
\caption{Wasserstein Error for Different $N$ at $t=0.5$}
\label{tab:sr-wasserstein-error}
\begin{tabular}{lccc}
\toprule
$N$ & \textbf{Mean $W_2$ Error} & \textbf{Standard Deviation} & \textbf{Theoretical
    Rate $N^{-1/2}$} \\
\midrule
50 & 0.0432 & 0.0081 & 0.0447 \\
100 & 0.0312 & 0.0057 & 0.0316 \\
200 & 0.0221 & 0.0040 & 0.0224 \\
\bottomrule
\end{tabular}
\end{table}


\section{Conclusion and Bridge to Alpha Generation}
\label{sec:sr-conclusion}

This chapter has applied the common-noise MFG framework to a concrete problem in
systemic risk and optimal bailout design. We derived the FBSDE system
characterising the equilibrium, solved the linear-quadratic case in closed form, and
quantified the $O(1/N)$ approximation error. The solution provides clear policy
implications regarding the trade-off between bailout generosity and moral hazard,
and the stabilising role of counter-cyclical intervention.

The systemic risk model serves as a warm-up for the empirical finance application
of Chapters 8--9. The key concepts --- distributional disequilibrium measured by the
Wasserstein distance, the impact of common noise on aggregate dynamics, and the
equilibrium response of individual agents --- carry over directly to the portfolio
choice setting. In Chapter 8, we replace banks with institutional investors,
reserves with portfolio positions, and bailouts with trading signals. The crowding
measure $c_t = W_2(\mu_t, \mu_t^*)$ will play the role of the Wasserstein distance
between the empirical and equilibrium measures, providing a theoretically grounded
signal for alpha generation.
